\documentclass{article}
\usepackage{amsmath,amssymb}
\usepackage[vocab={numbers,torus}]{leanify}
\newtheorem{lemma}{Lemma}
\begin{document}

\section{Conventions}

Throughout, $C$ and $C_m$ denote positive constants.
We write $N\varphi$ for the derivative of $\varphi$ in the direction $(\sqrt2-1,1)$.
We write $\|F\|_{C_y^m}$ for the $C^m$ norm of $F$.
We write $(a,b)_{\gcd}$ for the greatest common divisor of $a$ and $b$.

\section{Numbers}

\begin{lemma}\label{primes}
$7$ is prime and $9$ is a perfect square, and $9$ is not prime.
\end{lemma}

\begin{lemma}\label{even}
Let $n$ be an even integer. Then $n^2$ is even.
\end{lemma}

\begin{lemma}\label{gcd}
$(12,18)_{\gcd}=6$ and $\sqrt2$ is irrational.
\end{lemma}

\begin{lemma}\label{constants}
For every positive integer $m$ there is $c>0$ such that $c\,n^m\le n^{m+1}$ for every positive integer $n$.
\end{lemma}

\begin{lemma}\label{natural}
Let $n$ be a natural number. Then $n\ge0$.
\end{lemma}

\begin{lemma}\label{small}
For $0<h\ll\lambda\ll1$, $h<\lambda<1$.
\end{lemma}

\section{The torus}

\begin{lemma}\label{inverse}
Let $F$ be a smooth function on $\mathbb T^2$ of zero Haar mean. Then there is exactly one smooth function $\varphi$ on $\mathbb T^2$ of zero Haar mean with $N\varphi=F$. Moreover, $\|\varphi\|_{C_y^m}\le C_m\|F\|_{C_y^{m+4}}$ for every integer $m\ge0$ and every smooth function $\varphi$ on $\mathbb T^2$ of zero Haar mean with $N\varphi=F$.
\end{lemma}

\end{document}
